\clearpage
\section{Conditional 90\% $CL_s$ calibration}\label{sec:calibration}
\newcommand{\CalComplete}{456}
\newcommand{\CalSpectra}{12,080,640}
\newcommand{\ValidationSpectra}{1,368,000}


The calibration target is $CL_s=0.10$ throughout. This extension retains the
nominal correlated Gaussian profile and the fixed-GP-mean likelihood, and
replaces their asymptotic tail probabilities with full-spectrum Poisson
pseudoexperiments. The observed comparisons cover \CalComplete{} of the 456
requested dataset/mass coordinates. The selected results use \CalSpectra{} calibration
spectra and \ValidationSpectra{} separately seeded validation spectra.
The signal parameter in the combination is one shared $\epsilon^2$; each
constituent receives its own correctly normalized signal template.

This is a calibration conditional on the reviewed kernel coordinates, signal
resolution, support, masks and conversion factors. The GP posterior and
count-dependent training errors are refitted for every pseudoexperiment.
Kernel optimization and the earlier support/model-selection procedure are
not repeated. The 2016 numerical exception remains. These boundaries matter
when comparing this study with the historical closure results.

\subsection{What the earlier closure studies established}

The recent v4.9 studies separate practical recovery from exact centering.
At 65 MeV, v4.9.1 reported a native-10\% mean pull of $-0.246$, with 90\%
interval $[-0.417,-0.076]$; an independent continuation gave $-0.284$.
The $\pm0.5$ practical band, adopted after seeing the result, accepted this offset, while the recorded
zero-null screen still resolved it. The corresponding 1\% source scaled
by ten, with a different threshold truth and 40--300 MeV support, gave
$+0.139$ and pull width $1.140$. These are distinct statements about mean,
width, generating shape and support.

The v4.9.5 support study found no edge satisfying its original uniform
$|\langle\mathrm{pull}\rangle|<0.5$ requirement. A documented post-phase-1
amendment accepted generally smaller-than-0.75 means with a 1.25 gross-bias
guard. The frozen 36 MeV edge was confirmed on independent backgrounds;
its three 55 MeV means remained $-0.773$, $-0.841$ and $-0.960$. Its other
nine tested means, at 60--70 MeV, lay between $-0.110$ and $+0.325$.
This supports the practical recovery claim made there, rather than an
absence of all conditional offsets.

The scaled-data composite also changed truth and support at its two
lower-exposure 65 MeV replacement points. It is therefore not a controlled
luminosity-scaling experiment at that coordinate. Historical reference and
injected fits reoptimized the GP, and their injected yields used each
background toy's fitted reference error. Here the physical injected yield
is fixed across an ensemble. A mean pull with a random fitted denominator
and the bias $E[\hat A]-A_{\mathrm{true}}$ are different quantities.
The new validation records both the fitted-yield mean and its deviation in
units of one fixed reference error. A nonzero value describes conditional
estimator miscentering; it does not measure a bias in the observed data or
identify a unique cause.

\clearpage
\subsection{Generating truths and repeated analysis}

Each individual dataset has two declared generating intensities: its
mass-local observed GP mean extended over the complete support, and an
archived alternative shape. For 2015 the alternative is the saved
\texttt{fShiftSigPowTail} expected-count spectrum. For 2016 it is the archived
threshold-qualified blend. For native 2021 it is the v4.9.5
\texttt{fSigPowExpQ}-anchored logistic--Chebyshev6 residual stress spectrum.
Native bins are summed to the exact production edges, with no additional
exposure factor or post-crop renormalization.

These alternatives retain their source-fit limitations. In particular,
the native-2021 construction was motivated by threshold recovery; its
85--300 MeV tail deviance per degree of freedom was 6.220. It is not a
globally qualified pure \texttt{fSigPowExpQ} truth. The data-selected
reverse-injection smooth shape used in the original 71 MeV diagnostic
remains outside these two calibration truths. Its separate follow-up status
is reported in Sec.~\ref{sec:reverse_truth_validation}. A failure under that
shape should not be equated to a historical closure result obtained under
a different truth and fitting procedure.

For the combination, the two scenarios are all constituent local-GP
truths and all constituent archived stress truths. Their endpoint envelope
covers this declared finite set; it does not cover every mixed assignment
of constituent truths or arbitrary background misspecification.

Given a truth $t$, pseudoexperiments fluctuate every support bin as
\begin{equation}
 N_{di}\sim\operatorname{Pois}\!\left(b^{(t)}_{di}
       +\epsilon^2 s^{(1)}_{di}\right).
\end{equation}
The complete signal, including the tails outside the moving mask, enters
generation. Each spectrum is processed through log-count training, the
updated $\alpha_i=1/N_i$ prescription for positive counts (and the archived
$\alpha_i=1$ convention for zero counts), and the GP posterior prediction.
The resulting count mean and covariance enter the same two likelihoods
used for the observed comparison. Both methods receive identical spectra.
The count-to-$\epsilon^2$ factors are kept at their frozen release values.

\subsection{Empirical likelihood-ratio tails}

Let $T_\mu$ denote the bounded one-sided $\widetilde q_\mu$ statistic of
Sec.~2, with the physical lower bound implemented in the denominator as prescribed
by Ref.~\cite{cowan}. The denominator uses the null refit when the unbounded
best-fit signal is negative, and $T_\mu=0$ when that best fit exceeds $\mu$. For each generating truth the calibrated quantity is
\begin{equation}
 CL_s^{(t)}(\mu)=
 \frac{P_{\mu,t}(T_\mu\ge T_\mu^{\mathrm{obs}})}
      {P_{0,t}(T_\mu\ge T_\mu^{\mathrm{obs}})}.
\end{equation}
Both probabilities are upper tails of the same statistic, including ties.
The larger truth-specific upper endpoint defines the reported finite-set
envelope. This applies the modified-frequentist construction of
Ref.~\cite{readcls}; it does not rescale a Gaussian significance or replace
the statistic by a fitted-yield proxy.

\clearpage
\subsection{Efficient sampling and numerical checks}

Pseudoexperiments are generated in equal strata from full-spectrum Poisson
proposal means at several injected strengths. Additional proposals shift
both the window and sidebands toward the observed tail. An approximate
linear influence calculation chooses those proposals only. Every spectrum
is still fitted, and its exact generating probability supplies the weight.
If $g_k(n)$ is proposal $k$, then
\begin{equation}
 g(n)=\frac{1}{K}\sum_{k=1}^K g_k(n),\qquad
 \widehat p_{\mu,t}=\frac{1}{N}\sum_j
 \frac{f_{\mu,t}(n_j)}{g(n_j)}
 \mathbf{1}[T_\mu(n_j)\ge T_\mu^{\mathrm{obs}}].
\end{equation}
This is the mixture-reuse identity discussed by Berns~\cite{berns}, applied
here to $CL_s$ tails rather than Feldman--Cousins ordering. The full Poisson
density ratio includes every support bin. Independent direct-Poisson
validation samples do not use the proposal mixture. An analytic one-bin
Poisson test with 120,000 draws separately checks the weighting and
stratified variance calculation.

The initial run uses 256 spectra per proposal stratum and 500 independent
validation spectra per truth at each of 0, 2 and 5 reference sigmas. Tail
effective sample sizes, B and S+B normalization checks and stratified
Monte Carlo errors accompany every endpoint. The approximate 95\% MC
interval describes numerical uncertainty in a 90\% limit. It is not an
expected experimental band or a change to the requested confidence level.
The displayed envelope shading takes the componentwise maxima of the two
truth-specific approximate intervals; it is not a simultaneous confidence band.
A resolved endpoint requires both tail effective sample sizes above 100,
MC relative half-width at most 10\%, a numerical bracket at most 1.5\%,
and the normalization and monotonicity checks in the archived protocol.

The implementation exploits the smooth kernel through eigenfeatures at
relative cutoff $10^{-15}$. Very small nuisance-covariance eigenvalues
below $10^{-5}$ in Poisson-variance units can also be removed. At every
coordinate, dense-reference comparisons must keep the mean discrepancy
below $0.001$ predictive standard deviations, covariance discrepancy below
$0.001$ of its maximum diagonal, and absolute signed-root and test-statistic
differences below $0.001$. Failed approximation checks restore the dense
implementation. Observed predictions always use the original dense GP.
Separate scalar-versus-batch likelihood checks, complete fit-failure
retention and source hashes prevent an optimizer or approximation failure
from being counted as successful calibration.

\subsubsection{Independent sampling refinement}

The selected result set contains 106 refined coordinates, including 3 second-attempt coordinates. Refinement eligibility uses only censoring and Monte Carlo diagnostics. It does not use validation outcomes or select a background method by its observed limit.

New proposal centers are placed near unresolved endpoints and extend scans without a crossing. For full-spectrum truth $t$ and signal $g$ per reference error, the local spacing is at most $0.75[\sum_i g_i^2/(t_i+a g_i)]^{-1/2}$. Dense proposal centers are separate from the inversion grid. Every original guard node is retained.

Refined truths use fresh independent banks with 512 draws per proposal, or 1,024 on a second attempt. Unrefined truths regenerate their original 256-draw banks and verify their full-array hashes. Old draws are never reweighted under a mixture adapted after inspecting them. The same 500 independent validation spectra per cell are rescored and counted once.

Original resolution tests still apply. Sampling checks also require the endpoint, bracket and slope evaluations to fall within the refined ranges, with normalization standard error at most 0.05. Caps or failed checks leave an explicit unresolved result.

At 1 selected coordinates, a stricter nuisance-eigenvalue cutoff of $10^{-7}$ passes all 18 original proposal checks and the complete extended-range audit. The discrepancy tolerances and twelve-mode cap are unchanged. Original rejected approximation checks remain in the record; failure of the stricter candidate restores the dense calculation. Actual reference-fit or scalar/batch disagreement remains fatal.

\paragraph{Bounded-memory execution.} 72 selected mass coordinates, containing 4,017,408 calibration spectra, use chunks of 128 spectra with one BLAS thread per worker. A source-derived peak-array budget of at most 8\,GiB per worker includes a 512\,MiB runtime allowance. For 2 combined coordinates, the resource policy permits one worker with up to 6\,GiB alongside at most one 4\,GiB worker, with an aggregate budget of at most 10\,GiB. For 3 combined coordinates, the resource policy permits one worker with up to 8\,GiB alongside at most one 4\,GiB worker, with an aggregate budget of at most 12\,GiB. These policies require a fresh memory-pressure check before each launch. Full-spectrum Poisson weights and the bounded statistic are checked against the original expressions and unsplit/scalar fits, with signed-root agreement within $2\times10^{-5}$ and $q$ agreement within $10^{-4}$. The likelihood, proposal laws and inference target are unchanged. Validation reuses the same 500 holdout spectra per truth and injected strength; these counts are not added to previous validation ensembles.

\paragraph{Scalar reference initialization.} 1 fixed-background scalar reference fit required a bracketed score-root initializer across 1 selected mass coordinate. The original solver is tried first and its restart must satisfy the same $2\times10^{-7}$ score threshold and all signed-root/$q$ agreement gates. The production batch statistic, GP convention, science proposal arrays, complete science banks and validation seeds are unchanged. The original failure and its successful replay diagnostic are retained. The original 18 numerical-audit draws are unchanged; fresh extended numerical-audit draws use the new source identity and remain separate from calibration and validation. A later reporting-schema error was resolved by reconstructing completion metadata from saved results. All numerical result fields were preserved exactly, with no numerical reexecution.

\begin{table}[H]\centering\small
\begin{tabular}{lrrrrr}\toprule
Scope & Complete & Profile resolved & Fixed resolved & Raw ratio & Calibrated ratio\\\midrule
2015, 100\% & 72/72 & 70 & 72 & 0.422 & 1.446\\
2016, 100\% & 142/142 & 133 & 141 & 0.522 & 1.333\\
2021, 10\% & 201/201 & 201 & 200 & 0.607 & 1.199\\
All three & 41/41 & 41 & 41 & 0.595 & 1.346\\
\bottomrule\end{tabular}
\caption{Pointwise completion and Monte Carlo precision status. The last two columns give median fixed/profiled observed-limit ratios on the same finite completed mass subset. A smaller observed limit is not by itself a sensitivity or coverage statement. Resolved counts refer to the declared MC gates, not unconditional background-model qualification.}
\label{tab:calibration_summary}\end{table}

%%LIMIT_2015%%
%%LIMIT_2016%%
%%LIMIT_2021%%
%%LIMIT_COMBINED%%
\clearpage
\subsection{Dependence on the generating background}
\fig{0.99}{../figures/truth_dependence.pdf}{
Ratio of the archived-stress endpoint to the local-GP endpoint, retaining
the two likelihood methods separately. Values above one identify an envelope
controlled by the stress shape; values below one identify a local-GP-controlled
endpoint. Open circles flag limited MC precision in either component. Missing
or right-censored endpoints remain gaps rather than finite ratios. The all-three
comparison contains two joint generating scenarios, not all mixed assignments.
}{fig:cal_truth_dependence}

This dependence is propagated through the same reviewed kernel coordinates
and full-spectrum retraining procedure in each case. It quantifies the effect
of these declared generating alternatives on the observed endpoint; it does
not qualify either alternative as a complete description of the true background.
In particular, the source-fit limitations documented above remain relevant
when an archived stress shape controls the envelope.
\clearpage
\subsection{Local p-values}
\fig{0.99}{../figures/local_pvalues.pdf}{
Asymptotic and calibrated local p-values. The calibrated value is the larger
B-only tail probability from the two declared generating truths. Downward
triangles mark asymptotic values below the displayed floor;
complete values remain in the CSV ledger. The empirical convention includes
the bounded-statistic atom and assigns $p_0=1$ when the observed discovery
statistic is zero; the earlier asymptotic display instead maps $Z=0$ to
$p_0=0.5$. Open circles mark limited precision or an importance-weighted
estimate within three MC standard errors above one, displayed at one;
the unclipped estimate is retained in the ledger.
Neither curve includes a global trials correction.
}{fig:cal_pzero}
\begin{table}[H]\centering\small
\begin{tabular}{lrrrr}\toprule
 & \multicolumn{2}{c}{Gaussian profile} & \multicolumn{2}{c}{Fixed GP mean}\\
Scope & $m$ (MeV) & Calibrated $p_0$ & $m$ (MeV) & Calibrated $p_0$\\\midrule
2015, 100\% & 51 & $4.35\times10^{-4}$$^{\dagger}$ & 51 & $3.26\times10^{-4}$$^{\dagger}$\\
2016, 100\% & 90 & $4.98\times10^{-4}$ & 121 & $6.24\times10^{-5}$\\
2021, 10\% & 77 & $7.80\times10^{-3}$ & 77 & $1.60\times10^{-4}$\\
All three & 90 & $6.65\times10^{-2}$ & 90 & $2.04\times10^{-2}$\\
\bottomrule\end{tabular}
\caption{Minimum calibrated local $p_0$ among completed points. The calibration takes the larger tail probability from the two declared truths. A dagger marks limited MC precision. Mass selection is descriptive; these are not global p-values. The asymptotic minima are retained in Table~\ref{tab:pzero}.}
\end{table}


\clearpage
\subsection{Conditional mean response}
\fig{0.99}{../figures/bias_by_truth.pdf}{
Background-only mean fitted yield in units of a fixed reference standard error.
Each point uses 500 independently generated full spectra. The reference scale
is constant within that mass ensemble; this is not the historical mean pull,
whose denominator was a random fitted uncertainty. Solid and dashed curves
retain the two generating truths separately. These conditional offsets do
not identify a bias of the observed data or isolate kernel-selection effects.
Shading shows approximate pointwise 95\% Monte Carlo intervals for the mean,
using the saved sample variance and Student-$t$ quantile; it is not simultaneous.
}{fig:cal_bias}

The comparison across truths tests whether a proposed calibration depends
strongly on a particular generating spectrum. The local-GP control itself
need not reproduce its original window prediction after retraining: the
prediction is estimated from fluctuated sidebands through a regularized
smoother. The archived stress shapes can introduce additional interpolation
errors. Neither mechanism is removed by merely replacing a linear-Gaussian
background constraint with its positive log-GP counterpart.

\clearpage
\subsection{Interpolation offsets}
The saved proposal diagnostics provide a separate, approximate check of this
structure without another toy study. Let $b_t$ be a generating spectrum and
$\bar b_t$ the prediction after training on its unfluctuated sidebands.
The linearized signal offset is
\begin{equation}
 \delta A_{\mathrm{lin}}=
 \frac{w^T V^{-1}(b_t-\bar b_t)}{w^T V^{-1}w},\qquad
 V=\operatorname{diag}(\bar b_t)+C
\end{equation}
for the profiled fit, with $C=0$ for the fixed-mean fit. The ledger retains
this deterministic residual projection beside the actual ensemble mean.
Agreement would support an interpolation offset as a sufficient explanation
under that truth; disagreement can also reflect nonlinear training and
finite-count effects. This diagnostic neither changes the toy statistic nor
identifies the effect of reoptimizing the kernels.
\begin{table}[H]\centering\small
\begin{tabular}{llrrrr}\toprule
 & & \multicolumn{2}{c}{Gaussian profile} & \multicolumn{2}{c}{Fixed GP mean}\\
Scope & Truth & $|\langle\hat A\rangle|$ & $|\langle\hat A\rangle-\delta A_{\rm lin}|$ & $|\langle\hat A\rangle|$ & $|\langle\hat A\rangle-\delta A_{\rm lin}|$\\\midrule
2015, 100\% & Local GP & 0.287 & 0.021 & 0.475 & 0.045\\
2015, 100\% & Archived stress & 0.250 & 0.026 & 0.430 & 0.037\\
2016, 100\% & Local GP & 0.204 & 0.023 & 0.296 & 0.036\\
2016, 100\% & Archived stress & 1.157 & 0.031 & 1.839 & 0.037\\
2021, 10\% & Local GP & 0.138 & 0.026 & 0.193 & 0.027\\
2021, 10\% & Archived stress & 0.360 & 0.029 & 0.498 & 0.033\\
All three & Local GP & 0.316 & 0.030 & 0.520 & 0.034\\
All three & Archived stress & 5.361 & 0.035 & 6.217 & 0.036\\
\bottomrule\end{tabular}
\caption{Descriptive medians over completed background-only mass ensembles, in units of the fixed $\sigma_{\rm ref}$ at each mass. The first column for each method measures the absolute ensemble offset; the second measures its absolute difference from the linearized deterministic residual projection. These summaries retain the generating truths separately and do not pool toy counts or test coverage.}
\label{tab:bias_projection}\end{table}


\clearpage
\subsection{Independent validation}
\fig{0.99}{../figures/validation_exclusion.pdf}{
Independent five-reference-sigma injections. Every marker represents one
mass, generating truth and method, with 500 spectra in that cell. The axes
compare the fraction excluding the true injected yield before and after
calibration; the guides mark 0.10. Results are not pooled across masses.
The calibrated decision uses the declared two-truth envelope. This is a
conditional repeated-analysis test, not a measurement of physics sensitivity.
}{fig:cal_validation}
\begin{table}[H]\centering\small
\begin{tabular}{lrrr}\toprule
Validation family & Cells tested & Raw flags & Calibrated flags\\\midrule
True-yield exclusion & 3648 & 2039 & 0\\
Background-only local rejection & 1824 & 952 & 0\\
\bottomrule\end{tabular}
\caption{One-sided binomial excess-rate screens with Holm adjustment at 0.05. Flags count rejected cell-level null hypotheses, not rejected toys. The exclusion null is 0.10; the local-rejection null is 0.05. Tests are separate families. Counts during an incomplete run are provisional because the family is still growing. No cell is discarded.}
\end{table}
The completed suite contains 0 adjusted excess-exclusion flags and 0 adjusted excess-local-rejection flags. All flagged cells, intervals, and MC precision statuses remain in the accompanying ledgers.


\clearpage
Each positive injection is tested at its true fixed physical signal yield.
For the envelope, a validation spectrum is excluded only when its
$CL_s$ is below 0.10 for both calibration truths. Background-only spectra
separately test local rejection at $p_0<0.05$. Exact binomial intervals and
one-sided excess-rate tests accompany the observed fractions; Holm
adjustment controls the reported family of validation tests. Absence of a
rejection in this finite suite would support only the stated conditional
procedure, rather than uniform or truth-independent coverage.
With 500 spectra per cell and 3,648 exclusion tests, the first Holm step
requires at least 81 exclusions (16.2\%) to reject the 10\% null at the
family level. The analogous local-rejection screen, with 1,824 tests,
starts at 48 of 500 (9.6\%) against a 5\% null. Thus an absence of flags
cannot certify coverage to percent-level accuracy at every mass.

\subsection{Auxiliary check of the known 71 MeV failure}\label{sec:reverse_truth_validation}

The original reverse-truth five-reference-sigma injections excluded the true yield in 349/500 fixed-background fits and 145/500 Gaussian-profiled fits. These are the recorded uncalibrated asymptotic results. The reverse truth is absent from the two-truth calibration envelope, so the main validation does not establish that this failure has been corrected.

This auxiliary result reconstructs the selected 71 MeV calibration proposal arrays and their hashes, then evaluates both calibration and validation statistics with dense GP refits and the full conditioned covariance. It uses the original 1,500 spectra (500 at each of zero, two and five reference sigmas), regenerated from the original seeds and checked against the saved fitted yields and signed likelihood roots. The two methods remain paired; these spectra are reused and are not pooled with the main validation.

The bounded statistic is evaluated at the original physical injected yield, without a Wald reconstruction. Calibration still takes the larger tail result from the same two declared truths. The reverse truth is a known out-of-envelope diagnostic: this retrospective result does not establish truth-independent or global coverage.

\begin{table}[H]\centering\small

\begin{tabular}{llrrrr}\toprule

Test & Method & Raw & Calibrated & Tail-ready & MC-resolved\\\midrule

B-only local $5\%$ & Gaussian profile & 1 & 19 & 356 & 356\\

B-only local $5\%$ & Fixed mean & 38 & 25 & 373 & 373\\

$2\sigma_{\rm ref}$ exclusion & Gaussian profile & 100 & 0 & 440 & 440\\

$2\sigma_{\rm ref}$ exclusion & Fixed mean & 315 & 0 & 416 & 416\\

$5\sigma_{\rm ref}$ exclusion & Gaussian profile & 145 & 3 & 369 & 369\\

$5\sigma_{\rm ref}$ exclusion & Fixed mean & 349 & 2 & 364 & 364\\

\bottomrule\end{tabular}

\caption{Counts out of 500 in each original reverse-truth ensemble. Positive injections test $CL_s<0.10$; B-only spectra test local $p_0<0.05$. The last two columns qualify the calibrated decision: tail-ready requires both truth-specific tail effective sample sizes of at least 100 and finite errors; MC-resolved additionally places the envelope of pointwise $\pm1.96$-SE tail estimates wholly on one side of the threshold. Point estimates retain every toy, including limited decisions. Exact binomial intervals in the CSV condition on the chosen calibration bank and do not include its finite Monte Carlo uncertainty. No counts are pooled across methods, strengths or studies.}

\label{tab:reverse_truth_validation}\end{table}

The five-reference-sigma exclusion point estimates fall to 3/500 for the Gaussian profile and 2/500 for the fixed mean. These values lie well below the nominal 10\% rate and indicate conservative behavior under this truth. However, 682 of the 3,000 method--toy decisions have limited Monte Carlo precision; this diagnostic does not resolve coverage to percent-level accuracy.


\clearpage
\subsection{Interpretation of the two methods}

%%CALIBRATED_COMPARISON%%

The fixed-background likelihood can be used as a test statistic whose
sampling distribution is calibrated. Its narrow uncorrected error is not
retained as an assumption of known background in the generating ensemble.
An unbiased fitted yield is not a prerequisite for correct frequentist
coverage: a reproducible conditional offset can enter the simulated tail
distributions. This does not protect against an unmodeled offset or a
missing generating background. Extraction centering, coverage and expected
power therefore remain separate checks.
Whether it gives a smaller observed bound after calibration is an empirical
comparison; it cannot be inferred from its much narrower asymptotic curve.
The Gaussian-profiled likelihood uses the correlated uncertainty during
fitting as well as in the repeated-analysis study. A global trials factor
addresses the mass search, and cannot repair a local extraction offset or
misestimated uncertainty.

These observed scans do not define a method-selection rule. A final choice
should use independent closure and expected-power evidence. Selecting the
smaller observed limit at each mass would define a third procedure whose
selection must be included in its calibration.

\subsection{Release boundary and reproducibility}

The calibration derivative is under
\path{study_results/v4p9p13_calibration_20260905/}. The protocol, historical
review, frozen source contract, per-coordinate result traces, validation
ledgers and final summary identify the exact computation. Seeds are
separated by calibration/validation purpose, scope, mass and truth.
Calibration banks reuse their draws across tested strengths; validation
seeds additionally distinguish the injected strength.
Full arrays have recorded hashes; large banks are reproducible
from the saved seed rules. The original v4.9.13 report is preserved.

The results in this extension do not constitute calibration of the entire
historical hyperparameter-selection procedure. A paired truth-by-kernel-policy
study would be needed to separate the effects of reoptimization from the
conditional offsets shown here. A final physics-coverage claim would also
require justified background alternatives over the complete support,
appropriate treatment of remaining common systematics, and resolution of
the disclosed 2016 numerical status. These are separate from Monte Carlo
precision on the conditional limits reported in this note.
